\section*{Chapitre \ref{ch:b2:fourier} --- Séries de Fourier}

\begin{solution}{exo:b2:fourier:1}
L'imparité tue les $a_n$. Pour $n \geq 1$~:
\[
b_n = \frac{2}{\pi}\int_0^{\pi} \sin nt\,\dd t
= \frac{2}{\pi}\cdot\frac{1 - (-1)^n}{n}
= \begin{cases} \frac{4}{\pi n} & n \text{ impair},\\ 0 & n
\text{ pair}. \end{cases}
\]
Donc $S(f)(t) = \frac{4}{\pi}\sum_{k\geq0} \frac{\sin\bigl((2k+1)t
\bigr)}{2k+1}$. Dirichlet en $t = \frac\pi2$ (un point de continuité,
valeur $1$)~: $\sin\bigl((2k+1)\frac\pi2\bigr) = (-1)^k$, donnant
\[
1 = \frac4\pi \sum_{k\geq0}\frac{(-1)^k}{2k+1}
\quad\Longrightarrow\quad
\sum_{k\geq0}\frac{(-1)^k}{2k+1} = \frac{\pi}{4}
\quad\text{(Leibniz)}.
\]
Au saut $t = 0$~: la série somme vers $0 = \frac{f(0^+) +
f(0^-)}{2}$, comme Dirichlet le prescrit (chaque terme s'annule~:
cohérent).
\end{solution}

\begin{solution}{exo:b2:fourier:2}
La parité tue les $b_n$~; $a_0 = \frac{1}{\pi}\int_{-\pi}^\pi\abs
t\,\dd t = \pi$, et pour $n \geq 1$~:
\[
a_n = \frac{2}{\pi}\int_0^\pi t\cos nt\,\dd t
= \frac{2}{\pi}\cdot\frac{(-1)^n - 1}{n^2}
= \begin{cases} -\frac{4}{\pi n^2} & n \text{ impair},\\ 0 & n
\text{ pair}, \end{cases}
\]
(une intégration par parties). D'où
\[
\abs t = \frac{\pi}{2} - \frac{4}{\pi}\sum_{k\geq0}
\frac{\cos\bigl((2k+1)t\bigr)}{(2k+1)^2} ,
\]
avec convergence \emph{normale} ($\sum (2k+1)^{-2} < \infty$) ---
comme le \cref{thm:b2:fourier:parseval}~(1) le prédit pour cette
fonction \omterm{def:b2:metric:continuity}{continue} et $C^1$ par morceaux. En $t = 0$~:
\[
0 = \frac\pi2 - \frac4\pi\sum_{k\geq0}\frac{1}{(2k+1)^2}
\quad\Longrightarrow\quad
\sum_{k\geq0}\frac{1}{(2k+1)^2} = \frac{\pi^2}{8} .
\]
En scindant $\sum \frac{1}{n^2}$ en parties impaire et paire~: $S =
\frac{\pi^2}{8} + \frac S4$, donc $S = \frac{\pi^2}{6}$~: Bâle de
nouveau.
\end{solution}

\begin{solution}{exo:b2:fourier:3}
Parité~: $b_n = 0$~; $a_0 = \frac{1}{\pi}\int_{-\pi}^{\pi} t^2 =
\frac{2\pi^2}{3}$~; deux intégrations par parties donnent $a_n =
\frac{4(-1)^n}{n^2}$ ($n \geq 1$). D'où
\[
t^2 = \frac{\pi^2}{3} + 4\sum_{n\geq1}
\frac{(-1)^n}{n^2}\cos nt
\qquad (\abs t \leq \pi),
\]
\omterm{def:b2:funcseq:series}{normalement} convergente. En $t = 0$~: $0 = \frac{\pi^2}{3} +
4\sum\frac{(-1)^n}{n^2}$, c'est-à-dire $\sum_{n\geq1}
\frac{(-1)^{n+1}}{n^2} = \frac{\pi^2}{12}$. Parseval~:
\[
\frac{1}{2\pi}\int_{-\pi}^{\pi} t^4\,\dd t = \frac{\pi^4}{5}
= \frac{a_0^2}{4} + \frac12\sum_{n\geq1} a_n^2
= \frac{\pi^4}{9} + 8\sum_{n\geq1}\frac{1}{n^4} ,
\]
donc $\sum \frac{1}{n^4} = \frac18\bigl(\frac{\pi^4}{5} -
\frac{\pi^4}{9}\bigr) = \frac{\pi^4}{90}$.
\end{solution}

\begin{solution}{exo:b2:fourier:4}
Si de plus $f$ est $C^1$ par morceaux~: Parseval (démontré) donne
$\norm f_2^2 = \sum\abs{c_n}^2 = 0$, et la stricte positivité de
l'intégrale de la fonction \omterm{def:b2:metric:continuity}{continue} $\abs f^2$ force $f = 0$.

Pour $f$ seulement \omterm{def:b2:metric:continuity}{continue}, admettons le Parseval général~: même
démonstration d'une ligne. (Esquisse de la voie par densité~:
l'approximation trigonométrique de type Fejér/Weierstrass montre que
les polynômes trigonométriques sont $\norm\cdot_2$-denses parmi les
fonctions \omterm{def:b2:metric:continuity}{continues} périodiques~; puisque $f \perp$ tous, $\norm
f_2^2 = \langle f, f - P\rangle \leq \norm f_2\norm{f - P}_2$ pour
des approximants $P$, forçant $\norm f_2 = 0$.)
\end{solution}

\begin{solution}{exo:b2:fourier:5}
Parité~: $b_n = 0$~;
\[
a_n = \frac{2}{\pi}\int_0^\pi \cos(\alpha t)\cos(nt)\,\dd t
= \frac{2}{\pi}\cdot
\frac{(-1)^n\,\alpha\sin(\pi\alpha)}{\alpha^2 - n^2}
\]
(produit-somme, puis intégrer~; $a_0 =
\frac{2\sin(\pi\alpha)}{\pi\alpha}$). Dirichlet en $t = \pi$ (un
point de continuité du prolongement périodique, dont les valeurs
unilatérales coïncident par parité)~:
\[
\cos(\pi\alpha)
= \frac{\sin(\pi\alpha)}{\pi\alpha}
+ \sum_{n\geq1} \frac{2\alpha\sin(\pi\alpha)}{\pi(\alpha^2 -
n^2)}\,(-1)^n\cos(n\pi)
= \frac{\sin(\pi\alpha)}{\pi}\Bigl(\frac{1}{\alpha} +
\sum_{n\geq1}\frac{2\alpha}{\alpha^2 - n^2}\Bigr),
\]
en utilisant $(-1)^n\cos n\pi = 1$. En divisant par
$\sin(\pi\alpha)/\pi$ on obtient le développement de la cotangente.
\end{solution}

\begin{solution}{exo:b2:fourier:6}
En itérant $c_n(f') = \iu n\,c_n(f)$ ($k$ fois, intégration par
parties à travers les morceaux $C^{k}$ avec des valeurs de bord
concordantes)~: $c_n(f^{(k)}) = (\iu n)^k c_n(f)$. Les coefficients
de $f^{(k)}$, \omterm{def:b2:metric:continuity}{continue} par morceaux, sont bornés (et même $\to 0$,
Bessel)~:
\[
\abs{c_n(f)} = \frac{\abs{c_n(f^{(k)})}}{\abs n^k}
= O\bigl(\abs n^{-k}\bigr) .
\]
\end{solution}

\begin{solution}{exo:b2:fourier:7}
Parseval pour $f$ et pour $f'$ (les deux légitimes~: $f$ est $C^1$,
$f'$ \omterm{def:b2:metric:continuity}{continue} par morceaux --- et même \omterm{def:b2:metric:continuity}{continue})~:
\[
\frac{1}{2\pi}\int \abs f^2 = \sum_{n\neq0} \abs{c_n}^2
\quad (c_0 = 0 \text{ par l'hypothèse de moyenne nulle}),
\qquad
\frac{1}{2\pi}\int \abs{f'}^2 = \sum_{n\neq0} n^2\abs{c_n}^2 .
\]
Terme à terme, $n^2\abs{c_n}^2 \geq \abs{c_n}^2$ pour $\abs n \geq
1$~: l'inégalité s'ensuit. L'égalité force $(n^2 - 1)\abs{c_n}^2 = 0$
pour tout $n$, c'est-à-dire $c_n = 0$ pour $\abs n \geq 2$~: $f(t) =
c_1\eu^{\iu t} + c_{-1}\eu^{-\iu t} = a\cos t + b\sin t$ (forme
réelle)~; réciproquement de telles $f$ donnent l'égalité.
\end{solution}

\begin{solution}{exo:b2:fourier:8}
D'après l'\cref{exo:b2:fourier:1}, $S_{2N-1}(x) =
\frac4\pi\sum_{k=0}^{N-1}\frac{\sin((2k+1)x)}{2k+1}$. En $x_N =
\frac{\pi}{2N}$, avec $u_k = (2k+1)x_N = \frac{(2k+1)\pi}{2N}$~:
\[
S_{2N-1}(x_N) = \frac{4}{\pi}\sum_{k=0}^{N-1}\frac{\sin u_k}{2k+1}
= \frac{4}{\pi}\sum_{k=0}^{N-1}\frac{\sin u_k}{u_k}\cdot
\frac{u_k}{2k+1}
= \frac{2}{\pi}\sum_{k=0}^{N-1}\frac{\sin u_k}{u_k}\cdot
\frac{\pi}{N},
\]
puisque $\frac{u_k}{2k+1} = \frac{\pi}{2N}$. Les points $u_k$ sont
les milieux des $N$ sous-intervalles de $\intcc{0}{\pi}$ de longueur
$\frac{\pi}{N}$~: la somme est une somme de Riemann aux milieux de la
fonction \omterm{def:b2:metric:continuity}{continue} $u \mapsto \frac{\sin u}{u}$, donc \omterm{def:b2:integration:improper}{converge} vers
\[
\frac{2}{\pi}\int_0^\pi \frac{\sin u}{u}\,\dd u
\approx \frac{2}{\pi}\times 1.8519 \approx 1.179 .
\]
Les sommes partielles près du saut dépassent la valeur $1$ de
$\approx 18\%$ du demi-saut à jamais~: le phénomène de Gibbs,
quantifié.
\end{solution}

\begin{solution}{exo:b2:fourier:9}
À partir de $\cos 3t = 4\cos^3t - 3\cos t$~:
\[
\cos^3 t = \frac{3\cos t + \cos 3t}{4},
\qquad
\sin^2t\,\cos t = \cos t - \cos^3 t
= \frac{\cos t - \cos 3t}{4} .
\]
Chacun est un polynôme trigonométrique, donc égal à sa propre série
de Fourier (unicité des coefficients~: deux développements
différeraient d'un polynôme trigonométrique dont tous les
coefficients sont nuls). Pour $\cos^3t$~: $a_1 = \frac34$, $a_3 =
\frac14$, tous les autres $a_n$ et tous les $b_n$ nuls~; $c_{\pm1} =
\frac38$, $c_{\pm3} = \frac18$. Pour $\sin^2t\cos t$~: $a_1 =
\frac14$, $a_3 = -\frac14$~; $c_{\pm1} = \frac18$, $c_{\pm3} =
-\frac18$.
\end{solution}

\begin{solution}{exo:b2:fourier:10}
Calcul direct~:
\[
c_n = \frac{1}{2\pi}\int_{-\pi}^{\pi}\eu^{(a - \iu n)t}\dd t
= \frac{\eu^{(a-\iu n)\pi} - \eu^{-(a - \iu n)\pi}}
{2\pi(a - \iu n)}
= \frac{(-1)^n\sinh(a\pi)}{\pi(a - \iu n)} ,
\]
en utilisant $\eu^{\pm\iu n\pi} = (-1)^n$. En $t = \pi$ le
prolongement périodique saute de $\eu^{a\pi}$ à $\eu^{-a\pi}$~;
Dirichlet (sommes partielles \omterm{def:b2:quadratic:adjoint}{symétriques}) donne
\[
\cosh(a\pi) = \sum_{n\in\Z}(-1)^n c_n\,
= \frac{\sinh(a\pi)}{\pi}\Bigl(\frac1a +
\sum_{n\geq1}\Bigl(\frac{1}{a - \iu n} +
\frac{1}{a + \iu n}\Bigr)\Bigr)
= \frac{\sinh(a\pi)}{\pi}\Bigl(\frac1a +
\sum_{n\geq1}\frac{2a}{a^2 + n^2}\Bigr) ,
\]
les parties imaginaires des termes appariés se compensant. En
divisant par $\sinh(a\pi)$~:
\[
\coth(\pi a) = \frac{1}{\pi a} +
\sum_{n\geq1}\frac{2a}{\pi(a^2 + n^2)} ,
\]
le jumeau hyperbolique de l'\cref{exo:b2:fourier:5}.
\end{solution}

\begin{solution}{exo:b2:fourier:11}
Commutativité~: substituer $s = x - t$ et utiliser la périodicité de
l'intégrande. Pour $c_n(f * g)$, l'intégrande $(x, t) \mapsto
f(x-t)g(t)\eu^{-\iu nx}$ est \omterm{def:b2:metric:continuity}{continu}~; les deux intégrales itérées
coïncident (toutes deux, en tant que fonctions de la borne supérieure
de la variable extérieure, s'annulent au bord gauche et ont la même
dérivée --- continuité plus \cref{thm:b2:integration:continuity}
justifient la dérivation de l'intégrale itérée). D'où
\[
c_n(f*g) = \frac{1}{2\pi}\int_{-\pi}^{\pi} g(t)\,\eu^{-\iu nt}
\Bigl(\frac{1}{2\pi}\int_{-\pi}^{\pi}
f(x-t)\,\eu^{-\iu n(x-t)}\dd x\Bigr)\dd t
= c_n(f)\,c_n(g),
\]
l'intégrale interne valant $c_n(f)$ pour tout $t$ (substitution et
périodicité). Enfin le \cref{lem:b2:fourier:kernel} dit que
$S_N(f)(x) = \frac{1}{2\pi}\int f(x+u)D_N(u)\dd u$~; la substitution
$u \mapsto -t$ et la parité de $D_N$ en font $(f * D_N)(x)$.
\end{solution}

\begin{solution}{exo:b2:fourier:12}
\begin{enumerate}
  \item Avec $w = r\eu^{\iu t}$~:
        \[
        P_r(t) = 1 + 2\,\Re\frac{w}{1 - w}
        = \Re\frac{1 + w}{1 - w}
        = \frac{1 - \abs w^2}{\abs{1 - w}^2}
        = \frac{1 - r^2}{1 - 2r\cos t + r^2} > 0 .
        \]
        Moyenne $1$~: l'intégration terme à terme de la série
        \omterm{def:b2:funcseq:series}{normalement} convergente ne garde que $n = 0$.
  \item Pour $\delta \leq \abs t \leq \pi$~: $\cos t \leq
        \cos\delta$, donc $P_r(t) \leq \frac{1 - r^2}{1 -
        2r\cos\delta + r^2}$, dont le dénominateur tend vers $2 -
        2\cos\delta > 0$ tandis que le numérateur tend vers $0$~:
        \omterm{def:b2:funcseq:def}{convergence uniforme} vers $0$ en dehors de tout voisinage de
        $0$.
  \item L'intégration terme à terme (\omterm{def:b2:funcseq:series}{convergence normale} en $t$)
        donne $(f * P_r)(x) = \sum_n r^{\abs
        n}c_n(f)\eu^{\iu nx}$. L'argument d'identité approchée~:
        avec moyenne $1$ et positivité,
        \[
        \abs{(f*P_r)(x) - f(x)}
        \leq \frac{1}{2\pi}\int_{-\pi}^{\pi}
        \abs{f(x-t) - f(x)}\,P_r(t)\,\dd t ,
        \]
        scinder en $\abs t = \delta$~: au plus $\varepsilon$
        (Heine) plus $2\norm f_\infty\sup_{\delta\leq\abs
        t\leq\pi}P_r \to \varepsilon$~: \omterm{def:b2:funcseq:def}{convergence uniforme} quand
        $r \to 1^-$. C'est la sommation d'Abel de la série de
        Fourier --- le jumeau de Fourier de la théorie au bord du
        chapitre sur les séries entières.
\end{enumerate}
\end{solution}

\begin{solution}{pb:b2:fourier:1}
\textbf{1.} En moyennant le \cref{lem:b2:fourier:kernel} sur $n = 0,
\dots, N-1$ (linéarité de l'intégrale) on obtient $\sigma_N(f)(x) =
\frac{1}{2\pi}\int f(x+u)F_N(u)\dd u$~; chaque $D_n$ a pour moyenne
$1$, donc $F_N$ a pour moyenne $1$.

\textbf{2.} Avec $\eu^{\iu u} - 1 = 2\iu\,\eu^{\iu
u/2}\sin\frac u2$~:
\[
\sum_{n=0}^{N-1}\sin\Bigl(\Bigl(n + \frac12\Bigr)u\Bigr)
= \Im\Bigl[\eu^{\iu u/2}\,\frac{\eu^{\iu Nu} - 1}{\eu^{\iu u}
- 1}\Bigr]
= \Re\,\frac{1 - \eu^{\iu Nu}}{2\sin\frac u2}
= \frac{1 - \cos Nu}{2\sin\frac u2}
= \frac{\sin^2\frac{Nu}2}{\sin\frac u2} .
\]
En divisant par $N\sin\frac u2$~:
\[
F_N(u) = \frac1N\sum_{n=0}^{N-1}
\frac{\sin\bigl((n+\frac12)u\bigr)}{\sin\frac u2}
= \frac{1}{N}\,
\frac{\sin^2\frac{Nu}{2}}{\sin^2\frac u2} \geq 0 .
\]

\textbf{3.} Sur $\delta \leq \abs u \leq \pi$~: $\sin^2\frac u2
\geq \sin^2\frac\delta2$ et $\sin^2\frac{Nu}2 \leq 1$~: $F_N
\leq \frac{1}{N\sin^2(\delta/2)} \to 0$, \omterm{def:b2:funcseq:def}{uniformément} là.

\textbf{4.} Par la moyenne unité, $\sigma_Nf(x) - f(x) =
\frac{1}{2\pi}\int\bigl(f(x+u) - f(x)\bigr)F_N(u)\dd u$. Étant donné
$\varepsilon$, Heine fournit $\delta$ avec $\abs{f(x+u) - f(x)}
\leq \varepsilon$ pour $\abs u \leq \delta$, \omterm{def:b2:funcseq:def}{uniformément} en $x$.
Alors, en utilisant $F_N \geq 0$ et sa moyenne unité,
\[
\abs{\sigma_Nf(x) - f(x)}
\leq \varepsilon +
2\norm f_\infty\cdot\frac{1}{2\pi}
\int_{\delta\leq\abs u\leq\pi}F_N
\leq \varepsilon + \frac{2\norm
f_\infty}{N\sin^2\frac\delta2}
\leq 2\varepsilon
\]
pour $N$ grand, \omterm{def:b2:funcseq:def}{uniformément} en $x$~: le théorème de Fejér.

\textbf{5.} $F_N$ est paire de moyenne $1$~: chaque moitié
$\intcc{0}{\pi}$, $\intcc{-\pi}{0}$ porte une moyenne $\frac12$.
Alors
\[
\sigma_Nf(x) - \frac{f(x^+)+f(x^-)}{2}
= \frac{1}{2\pi}\int_0^\pi\bigl(f(x+u) -
f(x^+)\bigr)F_N\,\dd u
+ \frac{1}{2\pi}\int_{-\pi}^0\bigl(f(x+u) -
f(x^-)\bigr)F_N\,\dd u ;
\]
sur chaque moitié, scinder en $\abs u = \delta$ où la limite
unilatérale est $\varepsilon$-proche, et laisser la question 3 tuer
la partie lointaine~: les deux intégrales tendent vers $0$. La
majoration~: $F_N \geq 0$ donne $\abs{\sigma_Nf(x)} \leq
\frac{1}{2\pi}\int\abs{f(x+u)}F_N \leq \norm f_\infty$.

\textbf{6.} Chaque $\sigma_N(f)$ est un polynôme trigonométrique
(une moyenne des $S_n(f)$, $n < N$), et $\sigma_N(f) \to f$
\omterm{def:b2:funcseq:def}{uniformément}~: le théorème de Weierstrass trigonométrique.

\textbf{7.} $c_n(f) = 0$ pour tout $n$ rend chaque $S_n(f) = 0$,
donc chaque $\sigma_N(f) = 0$~; par Fejér, $f = \lim\sigma_Nf =
0$. En appliquant ceci à une différence~: les fonctions \omterm{def:b2:metric:continuity}{continues}
périodiques sont déterminées par leurs \omterm{def:b2:fourier:coefficients}{coefficients de Fourier}.

\textbf{8.} $S_Nf$ est la projection orthogonale de $f$ sur
$\mathcal T_N$ (\cref{prop:b2:fourier:bessel}), donc elle minimise
$\norm{f - P}_2$ sur $P \in \mathcal T_N$~; puisque $\sigma_Nf
\in \mathcal T_{N-1} \subseteq \mathcal T_N$~:
\[
\norm{f - S_Nf}_2 \leq \norm{f - \sigma_Nf}_2
\leq \norm{f - \sigma_Nf}_\infty \to 0
\]
(l'inégalité du milieu car la moyenne de $\abs\cdot^2$ est au plus
le sup au carré). Pythagore $\norm f_2^2 = \sum_{\abs
n\leq N}\abs{c_n}^2 + \norm{f - S_Nf}_2^2$ passe alors à la
limite~: Parseval pour toute fonction $f$ \omterm{def:b2:metric:continuity}{continue} périodique.

\textbf{9.} Soit $t_1, \dots, t_p$ les sauts de $f$ sur une
période, $M = \norm f_\infty$. Pour $\eta$ petit, on définit $g = f$
en dehors des intervalles $\intoo{t_j - \eta}{t_j + \eta}$ et par la
corde affine sur chacun de ces intervalles~: $g$ est \omterm{def:b2:metric:continuity}{continue},
périodique, $\norm g_\infty \leq M$, et
\[
\norm{f - g}_2^2 \leq \frac{1}{2\pi}\,p\cdot(2M)^2\cdot2\eta
\leq \varepsilon^2
\]
pour $\eta$ petit. Bessel fait de $S_N$ une contraction pour
$\norm\cdot_2$, donc
\[
\norm{f - S_Nf}_2
\leq \norm{f - g}_2 + \norm{g - S_Ng}_2 + \norm{S_N(g -
f)}_2
\leq 2\varepsilon + \norm{g - S_Ng}_2 ,
\]
et la question 8 donne $\limsup_N\norm{f - S_Nf}_2 \leq
2\varepsilon$ pour tout $\varepsilon$~: $\norm{f - S_Nf}_2 \to
0$, et Pythagore fournit Parseval pour toute fonction $f$ \omterm{def:b2:metric:continuity}{continue}
par morceaux~: l'« admis » du \cref{thm:b2:fourier:parseval} est
maintenant un théorème.

\textbf{10.} Le signal carré a $\norm f_\infty = 1$, donc la
question 5 donne $\abs{\sigma_N(f)} \leq 1$ partout et pour tout
$N$ --- aucun dépassement, jamais --- tandis que
l'\cref{exo:b2:fourier:8} montre $\sup_xS_{2N-1}(f)(x) \to
\approx 1.179$. La propriété responsable~: $F_N \geq 0$, donc
$\sigma_Nf(x)$ est une \emph{moyenne} pondérée de valeurs de $f$
et ne peut jamais sortir de $\intcc{\min f}{\max f}$~; $D_N$ prend
des valeurs négatives, donc $S_N$ le peut.

\textbf{11.} $\abs{\sin N\theta} \leq N\abs{\sin\theta}$ par
récurrence ($\abs{\sin(N{+}1)\theta} \leq
\abs{\sin N\theta}\abs{\cos\theta} +
\abs{\cos N\theta}\abs{\sin\theta} \leq
(N+1)\abs{\sin\theta}$)~: avec $\theta = \frac u2$,
\[
F_N(u) = \frac{\sin^2\frac{Nu}2}{N\sin^2\frac u2}
\leq \frac{N^2\sin^2\frac u2}{N\sin^2\frac u2} = N .
\]
La concavité de $\sin$ sur $\intcc{0}{\frac\pi2}$ donne
$\sin\frac u2 \geq \frac{u}{\pi}$ pour $0 \leq u \leq \pi$, donc
$F_N(u) \leq \frac{1}{N(u/\pi)^2} = \frac{\pi^2}{Nu^2}$.

\textbf{12.} Par parité et le scindage en $\frac1N$~:
\[
\frac{1}{2\pi}\int_{-\pi}^{\pi}\abs uF_N
= \frac1\pi\int_0^\pi uF_N
\leq \frac1\pi\Bigl(\int_0^{1/N}uN\,\dd u +
\int_{1/N}^{\pi}\frac{\pi^2}{Nu}\,\dd u\Bigr)
= \frac{1}{2\pi N} + \frac{\pi\ln(\pi N)}{N} .
\]
Pour $N \geq 2$~: $\ln(\pi N) \leq \bigl(1 +
\frac{\ln\pi}{\ln2}\bigr)\ln N \leq 2.66\ln N$ et
$\frac{1}{2\pi N} \leq \frac{\ln N}{N}$, donc le moment est
$\leq \frac{9\ln N}{N}$.

\textbf{13.} Pour $f$ $L$-lipschitzienne~:
\[
\abs{\sigma_Nf(x) - f(x)}
\leq \frac{1}{2\pi}\int\abs{f(x+u) - f(x)}F_N(u)\dd u
\leq L\cdot\frac{1}{2\pi}\int\abs uF_N
\leq \frac{9L\ln N}{N},
\]
\omterm{def:b2:funcseq:def}{uniformément} en $x$.

\textbf{14.} $c_0(e_1) = 0$ donne $S_0(e_1) = 0$, tandis que
$S_n(e_1) = e_1$ pour $n \geq 1$~: $\sigma_N(e_1) =
\frac{N-1}{N}e_1$, donc $\norm{\sigma_Ne_1 - e_1}_\infty =
\frac1N$. Même pour ce signal entier, à bande limitée, la vitesse
est $\frac1N$~: Fejér sature, exactement comme l'opérateur de
Bernstein sature à $\frac1n$ (Voronovskaya, dans le devoir du
chapitre sur les suites de fonctions).

\textbf{15.} Si $f$ s'annule sur $\intoo{x-\delta}{x+\delta}$,
alors $\sigma_Nf(x) = \frac{1}{2\pi}\int_{\delta \leq \abs u
\leq \pi}f(x+u)F_N(u)\dd u$, de valeur absolue au plus
$\frac{\norm f_\infty}{N\sin^2(\delta/2)} \to 0$~: les moyennes de
Cesàro en $x$ ne voient que $f$ près de $x$.

\textbf{16.} Étant donnés $\intcc ab$ et $\varepsilon$, choisir des
fonctions $\varphi^\pm$ \omterm{def:b2:metric:continuity}{continues} $1$-périodiques affines par
morceaux avec $\varphi^- \leq \mathbf 1_{\intcc ab} \leq \varphi^+$
et $\int_0^1(\varphi^+ - \varphi^-) \leq \varepsilon$ (trapèzes de
pentes sur des intervalles de longueur totale $\varepsilon$). Alors
\[
\limsup_N\frac{\#\{n \leq N : x_n \in \intcc ab\}}{N}
\leq \lim_N\frac1N\sum_{n\leq N}\varphi^+(x_n)
= \int_0^1\varphi^+ \leq b - a + \varepsilon ,
\]
et symétriquement $\liminf \geq b - a - \varepsilon$~: la
proportion tend vers $b - a$.

\textbf{17.} Pour $f = \eu^{2\iu\pi k\cdot}$ avec $k \neq 0$
l'hypothèse donne la limite $0 = \int_0^1f$~; pour $k = 0$ les deux
membres valent $1$~; la linéarité traite tout polynôme
trigonométrique. Pour $f$ \omterm{def:b2:metric:continuity}{continue} $1$-périodique et $\varepsilon >
0$, la question 6 (transportée par $t = 2\pi x$) fournit un polynôme
trigonométrique $P$ avec $\norm{f - P}_\infty \leq \varepsilon$~:
\[
\Bigl|\frac1N\sum_{n\leq N}f(x_n) - \int_0^1f\Bigr|
\leq 2\varepsilon +
\Bigl|\frac1N\sum_{n\leq N}P(x_n) - \int_0^1P\Bigr|
\longrightarrow 2\varepsilon .
\]
Avec la question 16~: $(x_n)$ est équirépartie.

\textbf{18.} Pour $k \neq 0$ et $\alpha$ irrationnel, $w =
\eu^{2\iu\pi k\alpha} \neq 1$~:
\[
\Bigl|\sum_{n=1}^{N}w^n\Bigr|
= \Bigl|\frac{w(w^N - 1)}{w - 1}\Bigr|
\leq \frac{2}{\abs{1 - w}} ,
\]
une borne indépendante de $N$~; en divisant par $N$ on obtient le
critère de Weyl, et la question 17 conclut~: $(\{n\alpha\})$ est
équirépartie.

\textbf{19.} Tout sous-intervalle reçoit une proportion asymptotique
égale à sa longueur, en particulier une infinité de points~:
$(\{n\alpha\})$ est dense. L'équirépartition est plus forte~: une
suite peut être dense tout en passant presque tout son temps dans un
coin (la densité dit où la suite \emph{va}, l'équirépartition dit
\emph{combien de fois}).

\textbf{20.} $2^n$ a pour premier chiffre $1$ si et seulement si
$10^m \leq 2^n < 2\cdot10^m$ pour un certain $m$, c'est-à-dire si et
seulement si $\{n\log_{10}2\} \in \intco{0}{\log_{10}2}$.
Irrationalité~: $\log_{10}2 = \frac pq$ donnerait $2^q = 10^p =
2^p5^p$, impossible pour $p \geq 1$ par unicité de la
factorisation. Le théorème de Weyl (question 18 avec $\alpha =
\log_{10}2$~; l'intervalle \omterm{def:b2:metric:topology}{semi-ouvert} est encadré entre des
intervalles fermés de longueurs voisines) donne la proportion
$\log_{10}2 \approx 0.301$~: les premiers chiffres de $2^n$ suivent
la loi de Benford.

\textbf{21.} $z$ est $C^1$, donc sa série de Fourier \omterm{def:b2:integration:improper}{converge}
\omterm{def:b2:funcseq:series}{normalement} de somme $z$ (\cref{thm:b2:fourier:parseval}~(1)), et
$c_n(z') = \iu n\,c_n$. Puisque $\abs{z'} = \frac{L}{2\pi}$ est
constant,
\[
\int_0^{2\pi}\abs{z'}^2\dd t
= 2\pi\Bigl(\frac{L}{2\pi}\Bigr)^{\!2}
= \frac{L^2}{2\pi} ,
\]
et Parseval appliqué à $z'$ \omterm{def:b2:metric:continuity}{continu} donne
$\frac{1}{2\pi}\int_0^{2\pi}\abs{z'}^2 =
\sum_n\abs{\iu nc_n}^2$, c'est-à-dire $\frac{L^2}{2\pi} = 2\pi\sum_n
n^2\abs{c_n}^2$.

\textbf{22.} Le Parseval polarisé $\frac{1}{2\pi}\int\conj
fg = \sum\conj{c_n(f)}c_n(g)$ découle de Parseval appliqué à
$f + g$ et $f + \iu g$ (\omterm{def:b2:quadratic:def}{identité de polarisation}), toutes deux
\omterm{def:b2:metric:continuity}{continues}. Avec $f = z$, $g = z'$~:
\[
A = \frac12\,\Im\int_0^{2\pi}\conj z\,z'
= \pi\,\Im\sum_n\conj{c_n}(\iu n c_n)
= \pi\sum_n n\abs{c_n}^2 .
\]

\textbf{23.} En combinant les questions 21--22~:
\[
L^2 - 4\pi A
= 4\pi^2\sum_n n^2\abs{c_n}^2 -
4\pi^2\sum_n n\abs{c_n}^2
= 4\pi^2\sum_{n\in\Z}(n^2 - n)\abs{c_n}^2 \geq 0 ,
\]
puisque $n^2 - n = n(n-1) \geq 0$ pour tout entier. L'égalité force
$c_n = 0$ pour tout $n \notin \{0, 1\}$~: $z(t) = c_0 +
c_1\eu^{\iu t}$, un cercle de centre $c_0$ et de rayon
$\abs{c_1} = \frac{L}{2\pi}$ (vitesse constante). Démonstration par
Hurwitz de l'inégalité isopérimétrique~: $A \leq \frac{L^2}{4\pi}$,
le cercle seul.

\textbf{24.} Cercle de rayon $R$~: $L = 2\pi R$, $A = \pi
R^2$~: $L^2 = 4\pi^2R^2 = 4\pi A$~: égalité. Carré de côté
$a$~: $L^2 = 16a^2 > 4\pi a^2 = 4\pi A$ (car $16 > 4\pi \approx
12.57$). Pour $n$ négatif, $n^2 - n = n(n - 1)$ est un produit
de deux entiers négatifs~: positif --- donc les modes tournant à
l'envers coûtent de l'aire deux fois. La vitesse constante est
entrée en question 21, convertissant $\int\abs{z'}^2$ en
$\frac{L^2}{2\pi}$~; pour une vitesse non constante,
Cauchy--Schwarz donne $\int\abs{z'}^2 \geq
\frac{(\int\abs{z'})^2}{2\pi} = \frac{L^2}{2\pi}$, donc
l'inégalité subsiste, le cercle restant le seul cas d'égalité.

\textbf{25.} (i) Tout découle de $F_N \geq 0$ (avec moyenne unité
et concentration)~; $D_N$ a moyenne unité et concentration
d'oscillation mais pas la positivité, et Gibbs en est le prix.
(ii) La sommabilité au sens de Cesàro de la série de Fourier
implique sa sommabilité au sens d'Abel avec la même somme
(Frobenius, démontré dans le devoir du chapitre sur les séries
entières) --- la voie du noyau de Poisson de
l'\cref{exo:b2:fourier:12} est exactement la méthode d'Abel.
(iii) Le théorème de Weyl n'a demandé que l'approximation uniforme
(question 6)~; l'inégalité isopérimétrique a demandé Parseval
lui-même (questions 8, 21--22). (iv) Le volume de Licence~3
démontre la complétude~: les exponentielles forment une base
hilbertienne de $L^2$, Parseval devient une isométrie d'espaces de
Hilbert, et le théorème de Fejér devient l'énoncé que cette
isométrie est calculable par moyennes positives.
\end{solution}
